\appendix

\section{Netlist do circuito completo}
\label{ap:netlist}

Netlist \texttt{rs2\_completo.cir} exatamente como foi simulada
(teste global com a referência ideal calculada pelo simulador):

{\small
\begin{verbatim}
* RS2 completo (21101175, modelo LM741_LEVEL1)
.include lm741_level1.lib
RA1 in1 suma2 11.1111k
RA2 in2 suma2 33.3333k
RAR vee suma2 7.5Meg
RA3 va suma2 100k
XA2 0 suma2 vcc vee vout1 LM741_LEVEL1
RAF2 suma2 vout1 100k
RA4 in3 suma1 10k
XA1 0 suma1 vcc vee va LM741_LEVEL1
RAF1 suma1 va 25k
VIN1 in1 0 SINE(0 20m 500)
VIN2 in2 0 SINE(0 20m 700)
VIN3 in3 0 SINE(0 20m 900)
RB3 in6 sumb2 10k
RB4 vb sumb2 90k
XB2 0 sumb2 vcc vee vout2 LM741_LEVEL1
RBF2 sumb2 vout2 90k
RB1 in4 sumb1 20k
RB2 in5 sumb1 10k
XB1 0 sumb1 vcc vee vb LM741_LEVEL1
RBF1 sumb1 vb 40k
VIN6 in6 0 SINE(0 20m 1500)
VIN4 in4 0 SINE(0 20m 1100)
VIN5 in5 0 SINE(0 20m 1300)
XC1 vout1 ng1 vcc vee na LM741_LEVEL1
XC2 vout2 ng2 vcc vee nb LM741_LEVEL1
XC3 np nn vcc vee out LM741_LEVEL1
RC1 ng1 na 10k
RCG ng1 ng2 5k
RC2 ng2 nb 10k
RC3 na nn 10k
RC4 nb np 10k
RC5 nn out 25k
RC6 np 0 25k
RL out 0 10k
VCC vcc 0 15
VEE vee 0 -15
BIDEAL ideal 0 V=112.5*V(in1)+37.5*V(in2)-31.25*V(in3)+25*V(in4)+50*V(in5)-112.5*V(in6)-2.5
.tran 0 20m 0 10u
.meas TRAN err_max MAX abs(V(out)-V(ideal))
.meas TRAN err_rms RMS (V(out)-V(ideal))
.meas TRAN vout_max MAX V(out)
.meas TRAN vout_min MIN V(out)
.end
\end{verbatim}
}

\section{Netlists dos blocos isolados}
\label{ap:blocos}

Bloco A (\texttt{rs2\_bloco\_a.cir}):

{\small
\begin{verbatim}
* RS2 bloco A: vout1 = -9*vin1 - 2.8*vin2 + 2.5*vin3 (21101175)
.include lm741_level1.lib
RA1 in1 suma2 11.1111k
RA2 in2 suma2 33.3333k
RAR vee suma2 7.5Meg
RA3 va suma2 100k
XA2 0 suma2 vcc vee vout1 LM741_LEVEL1
RAF2 suma2 vout1 100k
RA4 in3 suma1 10k
XA1 0 suma1 vcc vee va LM741_LEVEL1
RAF1 suma1 va 25k
VIN1 in1 0 SINE(0 0.1 300)
VIN2 in2 0 SINE(0 0.2 700)
VIN3 in3 0 SINE(0 0.4 1300)
RL vout1 0 10k
VCC vcc 0 15
VEE vee 0 -15
BIDEAL1 ideal1 0 V=-9*V(in1)-3*V(in2)+2.5*V(in3)+0.2
.tran 0 20m 0 10u
.meas TRAN vout1_max MAX V(vout1)
.meas TRAN err1_max MAX abs(V(vout1)-V(ideal1))
.end
\end{verbatim}
}

Bloco B (\texttt{rs2\_bloco\_b.cir}):

{\small
\begin{verbatim}
* RS2 bloco B: vout2 = 2*vin4 + 4*vin5 - 9*vin6 (21101175)
.include lm741_level1.lib
RB3 in6 sumb2 10k
RB4 vb sumb2 90k
XB2 0 sumb2 vcc vee vout2 LM741_LEVEL1
RBF2 sumb2 vout2 90k
RB1 in4 sumb1 20k
RB2 in5 sumb1 10k
XB1 0 sumb1 vcc vee vb LM741_LEVEL1
RBF1 sumb1 vb 40k
VIN6 in6 0 SINE(0 0.1 1300)
VIN4 in4 0 SINE(0 0.2 300)
VIN5 in5 0 SINE(0 0.2 700)
RL vout2 0 10k
VCC vcc 0 15
VEE vee 0 -15
BIDEAL2 ideal2 0 V=2*V(in4)+4*V(in5)-9*V(in6)
.tran 0 20m 0 10u
.meas TRAN vout2_max MAX V(vout2)
.meas TRAN err2_max MAX abs(V(vout2)-V(ideal2))
.end
\end{verbatim}
}

Bloco C (\texttt{rs2\_bloco\_c.cir}):

{\small
\begin{verbatim}
* RS2 bloco C: vout = 12.5*(vout2 - vout1), G=25 (21101175)
.include lm741_level1.lib
XC1 vout1 ng1 vcc vee na LM741_LEVEL1
XC2 vout2 ng2 vcc vee nb LM741_LEVEL1
XC3 np nn vcc vee out LM741_LEVEL1
RC1 ng1 na 10k
RCG ng1 ng2 5k
RC2 ng2 nb 10k
RC3 na nn 10k
RC4 nb np 10k
RC5 nn out 25k
RC6 np 0 25k
VN vout1 0 1
VP vout2 0 SINE(1 0.4 1k)
RL out 0 10k
VCC vcc 0 15
VEE vee 0 -15
.tran 0 5m 0 5u
.meas TRAN gain_diff FIND V(out)/0.4 AT=2.25m
.meas TRAN vout_avg AVG V(out)
.end
\end{verbatim}
}

\section{Variantes de ensaio}
\label{ap:variantes}

As seis netlists de resposta individual usam o mesmo circuito completo
do Apêndice~A, trocando só as fontes e as medidas. Exemplo para
$v_{in1}$ (\texttt{rs2\_vin1.cir}); nas demais muda apenas qual
entrada recebe o seno:

{\small
\begin{verbatim}
VIN1 in1 0 SINE(0 50m 1k)
VIN2 in2 0 0
VIN3 in3 0 0
VIN4 in4 0 0
VIN5 in5 0 0
VIN6 in6 0 0
.tran 0 5m 0 5u
.meas TRAN vout_pk FIND V(out) AT=2.25m
.meas TRAN vout_tr FIND V(out) AT=2.75m
.meas TRAN g_out_vin1 PARAM (vout_pk-vout_tr)/0.1
.meas TRAN v1_pk FIND V(vout1) AT=2.25m
.meas TRAN v1_tr FIND V(vout1) AT=2.75m
.meas TRAN g_vout1_vin1 PARAM (v1_pk-v1_tr)/0.1
.meas TRAN v2_pk FIND V(vout2) AT=2.25m
.meas TRAN v2_tr FIND V(vout2) AT=2.75m
.meas TRAN g_vout2_vin1 PARAM (v2_pk-v2_tr)/0.1
.meas TRAN vout_dc AVG V(out)
\end{verbatim}
}

A variante com o macromodelo comercial
(\texttt{rs2\_completo\_vendor.cir}) também usa o circuito do
Apêndice~A, com estas diferenças:

{\small
\begin{verbatim}
.include lm741_vendor.lib          (macromodelo comercial da TI)
XA1 ... XC3 usam o subcircuito LM741 no lugar de LM741_LEVEL1
.options method=Gear
.options cshunt=10p abstol=1e-9 vntol=10u
.tran 0 10m 0 10u
\end{verbatim}
}

Essas opções numéricas foram necessárias porque o macromodelo
comercial tem dificuldade de convergência com a configuração padrão do
simulador; 10 ms cobrem um período completo do teste (as frequências
são múltiplas de \SI{100}{\hertz}).

\section{Código gerador}
\label{ap:codigo}

Esquemáticos e netlists são gerados por
\texttt{experiments/rs2-somador-subtrator/rs2\_somador\_subtrator.py}
(no repositório), executado com \texttt{./claw-spice code build}. As
simulações rodam uma por contêiner:
\texttt{./claw-spice sim run <netlist> -{}-timeout 900}, em paralelo.
O modelo \texttt{LM741\_LEVEL1} usado em todas as netlists está na
Seção~\ref{sec:modelo}.
